%% Chapter 5 -----------------------------------------------------------------
\chapter{Opening and Combining Data}
\label{ch:opening}

This chapter describes the file types the main window reads, the ways of opening them,
and the rules and options for combining several files into one dynamic spectrum.

\section{Supported files}
\index{FITS file!supported formats}

The main window reads CALLISTO dynamic spectra in FITS format, compressed or
uncompressed:
\begin{itemize}
  \item uncompressed files: \file{.fit}, \file{.fits};
  \item gzip-compressed files: \file{.fit.gz}, \file{.fits.gz}.
\end{itemize}
Upper-case suffixes such as \file{.FITS} are accepted. Combination relies on the standard
CALLISTO file name, \file{STATION_YYYYMMDD_HHMMSS_FC}, from which the station, start time
and focus code are read. Files converted from the Learmonth spectrograph
(Section~\ref{sec:learmonth}) follow the same conventions.

\subsection{ARTEMIS-IV files}
\label{sec:artemis}
\index{ARTEMIS-IV}

FITS files from the ARTEMIS-IV radio spectrograph (Thermopylae, Greece) in its
\texttt{ARTLOOK} format are recognized automatically. On loading:
\begin{itemize}
  \item the hours-based UT time axis is converted to the internal time axis;
  \item intensities are labeled and scaled as ADC counts, with the instrument's
    conversion of 58.51~counts per dB used for the dB scale; and
  \item the observation opens background subtracted, because the receiver's channel gains
    span more than an order of magnitude.
\end{itemize}

\section{Opening files}
\label{sec:opening-files}
\index{opening files}

Use any of the following:
\begin{itemize}
  \item \menu{File > Open}, the \gui{Open / Load}\index{toolbar!Open / Load} toolbar tool, or \kbd{Ctrl+O};
  \item drag the files onto the main window;
  \item \menu{File > Recent Files}\index{recent files} (Section~\ref{sec:recent-files});
  \item the \gui{Import} button of a downloader (Chapters~\ref{ch:downloader}
    and~\ref{ch:other-radio}).
\end{itemize}

When one file is selected, it is displayed at once. When several files are selected, the
application determines whether they can be combined in time, in frequency, or in both
(Section~\ref{sec:combine-rules}) and, if so, combines them into a single spectrum. If the
selection cannot be combined, a message explains why
(Figure~\ref{fig:combine-error}). If the current work has unsaved project changes, you are
asked whether to save them first.

After loading, the analysis controls and the \gui{Graph Properties} panel become active,
and the sidebar's \gui{Timeline} panel lists the observations that make up the dataset.

\subsection{Recent files}
\label{sec:recent-files}
\index{recent files}

\menu{File > Recent Files} lists the last ten FITS selections. A combined set is a single
entry, combined again when reopened; hovering over an entry shows its files. An entry
whose files have been moved or deleted is removed from the list with a message.
\gui{Clear Recent Files} empties the list.

\section{Combining files}
\label{sec:combine-rules}
\index{combining files}

CALLISTO observations are fragmented in time (15-minute files) and often in frequency
(several receivers or focus codes). The application can merge them along either axis or
both. The rules are summarized in Table~\ref{tab:combine-rules}.

\begin{table}[!htbp]
  \centering
  \caption{Requirements for combining FITS files.}
  \label{tab:combine-rules}
  \small
  \begin{tabularx}{\linewidth}{@{}P{0.2\linewidth}L@{}}
    \toprule
    \thead{Combination} & \thead{Requirements} \\
    \midrule
    Time & Same station and focus code; consecutive observations whose start times are 750--1050~s apart (normally 15~minutes). Sequences may continue across UTC midnight. \\
    Frequency & Same station, date and observation start time; distinct focus codes; identical time axes. Bands may be separated by a gap or may overlap (Section~\ref{sec:gap-overlap}). \\
    Time~\(\times\)~frequency & A complete grid of at least two consecutive timestamps and at least two focus codes, with the same set of focus codes at every timestamp. \\
    \bottomrule
  \end{tabularx}
\end{table}

A selection that contains extra files is searched for a combinable subset; for example,
selecting a whole day of files together with one stray file still finds the grid within
it. All files of one combined dataset must come from the same station.

For a time~\(\times\)~frequency grid, the bands are first combined at each timestamp with
one shared gap and overlap policy; the application checks that every result has the same
frequency grid, and then joins the results in time order. Incomplete grids, duplicate
timestamp and focus-code pairs, non-consecutive timestamps and incompatible axes are
rejected with a detailed message.

\screenshot[0.62\linewidth]{downloader/time_combine_err}{Error message for files that cannot be combined.}{fig:combine-error}{Warning that names the offending timestamps.}

Combined datasets keep their provenance: exported FITS files, projects and reports record
the original list of source files and the combination method (\texttt{time},
\texttt{frequency} or \texttt{time\_frequency}).

\subsection{Frequency gaps and overlaps}
\label{sec:gap-overlap}
\index{frequency gap}\index{frequency overlap}

When the selected bands leave a gap between them, or overlap, the \gui{Frequency Combine
Options} dialog asks how to handle it before the combined spectrum is imported
(Figures~\ref{fig:gap-options} and~\ref{fig:overlap-options}). The detected gap or overlap
range is shown at the top of the dialog. The chosen option is stored in the metadata of the
combined dataset.

\screenshot[0.62\linewidth]{downloader/freq_combine_gap_options}{Gap handling in the \gui{Frequency Combine Options} dialog.}{fig:gap-options}{Frequency Combine Options dialog for a gap.}

\begin{description}
  \item[Interpolated background fill] fills the gap with background values interpolated
    between the edges of the two bands.
  \item[Gray-hatched blank gap] leaves the gap without data and displays it as a
    gray-hatched region, so that it is clearly distinguished from observed data.
  \item[Average edge background fill] fills the gap with the average background level of
    the band edges on either side.
  \item[Zero fill] fills the gap with zeros.
\end{description}

\screenshot[0.62\linewidth]{downloader/freq_combine_overlap_options}{Overlap handling with a connection frequency.}{fig:overlap-options}{Frequency Combine Options dialog for an overlap.}

\begin{description}
  \item[Split at connection frequency] uses the lower band below the
    \gui{Connection frequency}\index{frequency overlap!connection frequency} and the higher band above it. The connection frequency is
    proposed automatically and can be edited.
  \item[Use lower-frequency file] keeps the data of the lower-frequency file in the
    overlapping region. The file name is shown in the option when there is a single overlap.
  \item[Use higher-frequency file] keeps the data of the higher-frequency file in the
    overlapping region.
  \item[Reject overlap] cancels the combination.
\end{description}

\section{The Combine FITS Files dialog}
\label{sec:combine-dialog}

\menu{File > Combine FITS Files...} opens a dialog for selecting, validating and previewing
a combination before it is imported (Figure~\ref{fig:combine-dialog}).

\begin{steps}
  \item Click \gui{Select FITS Files} and choose the files. The dialog inspects them and
    reports whether they form a valid time, frequency or time~\(\times\)~frequency
    combination.
  \item For a frequency combination, choose the gap and overlap handling and, if needed,
    the connection frequency in the \gui{Frequency Combine Options} group. In this dialog
    the overlap choices are labeled \gui{Keep low band in overlap} and \gui{Keep high
    band in overlap}.
  \item Click \gui{Combine and Preview}\index{Combine FITS Files dialog} to see the combined spectrum in the dialog.
  \item Click \gui{Import to Analyzer} to load the result into the main window.
\end{steps}

\screenshot[0.6\linewidth]{dialogs/combine_fits_dialog}{The \gui{Combine FITS Files} dialog.}{fig:combine-dialog}{The Combine FITS Files dialog after Select FITS Files and Combine and Preview, showing the gap/overlap options group and the combined preview image.}

\section{Viewing the FITS header}
\label{sec:fits-header}
\index{FITS header}

\menu{View > View FITS Header} shows the primary header of the loaded file
(Figure~\ref{fig:fits-header}). \gui{Save as .txt} writes the header to a text file.

\screenshot[0.62\linewidth]{dialogs/fits_header_viewer}{The FITS header viewer.}{fig:fits-header}{The FITS header viewer for a loaded CALLISTO file, showing the header cards and the Save as .txt and Close buttons.}
